\documentclass{article}
\usepackage{graphicx} % Required for inserting images
\usepackage{amsmath}
\title{Solutions to Worksheet \#2}
\author{}
\date{August 3, 2026}

\begin{document}

\maketitle
\begin{enumerate}
\item \textbf{Solution.}
For the change of basis from $e$ to $m$, the number
$M_{\lambda,\mu}(e,m)$ counts fillings of the rows of shape $\lambda$
with content $\mu$ such that every row is \emph{strictly increasing}.
There is no condition down the columns.  Rows in different positions
are regarded as different, even when they have the same length.

Here the shape $(2^2,1^4)=(2,2,1,1,1,1)$ has two rows of length $2$
and four one-cell rows.  Call the long rows $A$ and $B$.  The content
$(3^2,1^2)=(3,3,1,1)$ means that the complete filling must contain
three $1$'s, three $2$'s, one $3$, and one $4$.  Because a long row is
strictly increasing, it cannot contain two equal entries.

We count according to where the exceptional entries $3$ and $4$ occur.

\begin{enumerate}
\item Suppose neither $3$ nor $4$ is in a long row.  Then both long
rows must be $[1,2]$.  The four one-cell rows contain $1,2,3,4$ in any
order, giving
\[
4!=24
\]
fillings.

\item Suppose exactly one of $3,4$ is in a long row.  Choose which
exceptional entry is used ($2$ choices), choose whether it is in $A$
or $B$ ($2$ choices), and choose whether its companion is $1$ or $2$
($2$ choices).  The other long row is then forced to be $[1,2]$.
The four one-cell rows contain two equal entries and two other entries,
so they can be filled in
\[
\frac{4!}{2!}=12
\]
ways.  This case contributes
\[
2\cdot2\cdot2\cdot12=96.
\]

\item Suppose both $3$ and $4$ occur in the long rows.  There are
three possibilities.
\begin{enumerate}
\item If $3$ and $4$ are together, choose their long row in $2$ ways.
The other long row must be $[1,2]$, and the one-cell rows contain
$1,1,2,2$.  This gives
\[
2\cdot\frac{4!}{2!2!}=12.
\]

\item If $3$ and $4$ are in different long rows and have the same
companion, choose that companion in $2$ ways and choose which long row
contains $3$ in $2$ ways.  The one-cell rows contain three copies of
one of $1,2$ and one copy of the other, so this gives
\[
2\cdot2\cdot\frac{4!}{3!}=16.
\]

\item If $3$ and $4$ are in different long rows and have different
companions, choose which long row contains $3$ in $2$ ways and choose
which exceptional entry is paired with $1$ in $2$ ways.  The one-cell
rows contain $1,1,2,2$, so this gives
\[
2\cdot2\cdot\frac{4!}{2!2!}=24.
\]
\end{enumerate}
Thus this case contributes $12+16+24=52$.
\end{enumerate}

Adding the three cases, we obtain
\[
\boxed{M_{(2^2,1^4),(3^2,1^2)}(e,m)=24+96+52=172.}
\]

\item \textbf{Solution.}
We use row-strict tableaux of shape $\lambda$ and content $1^n$.
The content $1^n=(1,1,\ldots,1)$ means that the entries
$1,2,\ldots,n$ each appear exactly once.

First choose which $\lambda_1$ entries go in the first row, then which
$\lambda_2$ of the remaining entries go in the second row, and so on.
Once the entries of a row have been chosen, there is only one legal way
to place them: in increasing order.  Therefore
\[
\boxed{M_{\lambda,1^n}(e,m)
=\binom{n}{\lambda_1,\ldots,\lambda_\ell}
=\frac{n!}{\lambda_1!\cdots\lambda_\ell!}.}
\]

\item \textbf{Solution.}
The content $(n)$ means that every one of the $n$ cells must contain
$1$.  A strictly increasing row cannot contain two $1$'s.  Therefore
every row must have length $1$, which is exactly the shape
$\lambda=1^n$.  For this shape there is one filling: put $1$ in every
cell.  Hence
\[
\boxed{M_{\lambda,n}(e,m)=
\begin{cases}
1,&\lambda=1^n,\\
0,&\text{otherwise}.
\end{cases}}
\]

\item \textbf{Solution.}
Assume $n\geq2$.  The content $(n-1,1)$ means that the tableau has
$n-1$ entries equal to $1$ and one entry equal to $2$.

Since rows are strictly increasing, no row can contain two $1$'s.
Thus every row has length $1$, except possibly one row of length $2$.
Consequently, the only possible shapes are $1^n$ and
$(2,1^{n-2})$.

For shape $1^n$, choose which of the $n$ one-cell rows contains the
unique $2$; all other rows contain $1$.  This gives $n$ fillings.  For
shape $(2,1^{n-2})$, the long row is forced to be $[1,2]$, and every
one-cell row is forced to contain $1$, giving one filling.  Therefore
\[
\boxed{M_{\lambda,(n-1,1)}(e,m)=
\begin{cases}
n,&\lambda=1^n,\\
1,&\lambda=(2,1^{n-2}),\\
0,&\text{otherwise}.
\end{cases}}
\]

\item \textbf{Solution.}
The statement is false.  Take
\[
\lambda=(4,1,1),\qquad \mu=(3,3).
\]
Both are partitions of $6$, and $\lambda>_{\mathrm{lex}}\mu$ because
$4>3$.  Their conjugates are
\[
\lambda'=(3,1,1,1),\qquad \mu'=(2,2,2).
\]
Now $\mu'<_{\mathrm{lex}}\lambda'$, rather than
$\mu'>_{\mathrm{lex}}\lambda'$.  This is a counterexample.

\item \textbf{Solution.}
Let
\[
E(t)=\prod_{i\geq1}(1+x_it)=\sum_{r\geq0}e_rt^r
\]
and let $\omega$ be a primitive cube root of unity, so
$1+\omega+\omega^2=0$.  Since
\[
1+x+x^2=(1-\omega x)(1-\omega^2x),
\]
we obtain
\[
\boxed{\prod_{i\geq1}(1+x_i+x_i^2)
=E(-\omega)E(-\omega^2)
=\left(\sum_{r\geq0}(-\omega)^re_r\right)
 \left(\sum_{s\geq0}(-\omega^2)^se_s\right).}
\]
This expression has integer coefficients because the two factors are
complex conjugates.  For example, its first homogeneous pieces are
\[
1+e_1+(e_1^2-e_2)+(e_1e_2-2e_3)+\cdots.
\]

\item \textbf{Solution.}
\begin{enumerate}
\item Only the first five rows and first three columns can be nonzero.
The row sums are $(3,1,1,1,1)$ and the column sums are $(4,2,1)$.
The row of sum $3$ must be $(1,1,1)$.  After removing it, the four
remaining rows each contain one $1$, and the remaining column sums are
$(3,1,0)$.  Choose which three of the four rows place their $1$ in the
first column; the remaining row uses the second column.  Thus
\[
\boxed{M_{(4,2,1),(3,1^4)}(e,m)=\binom43=4.}
\]

\item Transposition is a bijection.  A matrix with row sums $\mu_i$
and column sums $\lambda_j$ is sent by transposition to a matrix with
row sums $\lambda_j$ and column sums $\mu_i$.  Therefore
\[
\boxed{M_{\lambda,\mu}(e,m)=M_{\mu,\lambda}(e,m).}
\]

\item In the tableaux interpretation, row $j$ of the Young diagram of
$\lambda$ is filled with distinct positive integers, written in
increasing order, and integer $i$ occurs $\mu_i$ times.  Define
$a_{ij}=1$ when integer $i$ occurs in tableau row $j$, and define
$a_{ij}=0$ otherwise.  Distinctness within a tableau row makes this a
$\{0,1\}$-matrix.  Its row sum is $\mu_i$ and its column sum is
$\lambda_j$.

Conversely, from such a matrix, fill tableau row $j$ with the indices
$i$ for which $a_{ij}=1$, in increasing order.  These constructions
are inverse to one another, so they give the required bijection.
\end{enumerate}

\item \textbf{Solution.}
For the change of basis from $h$ to $m$, the number
$M_{\lambda,\mu}(h,m)$ counts fillings of the rows of shape $\lambda$
with content $\mu$ such that every row is \emph{weakly increasing}.
Again, there is no condition down the columns, and equal-length rows in
different positions are distinguished.

We have two long rows $A,B$ of length $2$ and four one-cell rows.  We
must place three $1$'s, three $2$'s, one $3$, and one $4$.  The main
difference from Exercise 1 is that repeated entries are now allowed in
a long row: for example, $[1,1]$ and $[2,2]$ are legal.

We count according to how many of $3,4$ occur in the long rows.

\begin{enumerate}
\item Suppose neither $3$ nor $4$ is in a long row.  Each of $A,B$ is
one of
\[
[1,1],\qquad [1,2],\qquad [2,2].
\]
The pair $([1,1],[1,1])$ is impossible because it uses four $1$'s,
and $([2,2],[2,2])$ is impossible for the same reason with $2$'s.
The legal possibilities contribute as follows:
\[
\begin{array}{c|c|c|c}
\text{types of }A,B&\text{orders}&\text{entries in one-cell rows}
&\text{count}\\ \hline
\{[1,1],[1,2]\}&2&2,2,3,4&2\cdot 4!/2!=24\\
\{[1,1],[2,2]\}&2&1,2,3,4&2\cdot 4!=48\\
\{[1,2],[1,2]\}&1&1,2,3,4&4!=24\\
\{[1,2],[2,2]\}&2&1,1,3,4&2\cdot 4!/2!=24
\end{array}
\]
Thus this case contributes $120$.

\item Suppose exactly one of $3,4$ is in a long row.  Choose that
exceptional entry in $2$ ways and choose its long row in $2$ ways.
Its companion is either $1$ or $2$.

If its companion is $1$, the other long row may be $[1,1]$, $[1,2]$,
or $[2,2]$.  The corresponding one-cell rows can be filled in
\[
\frac{4!}{3!}=4,qquad
\frac{4!}{2!}=12,qquad
\frac{4!}{2!}=12
\]
ways, respectively.  This gives $4+12+12=28$ fillings.  If the
companion is $2$, the same argument with $1$ and $2$ interchanged gives
another $28$.  Therefore this case contributes
\[
2\cdot2\cdot(28+28)=224.
\]

\item Suppose both $3$ and $4$ are in the long rows.
\begin{enumerate}
\item They may be together in one long row, which can be chosen in
$2$ ways.  The other long row may be $[1,1]$, $[1,2]$, or $[2,2]$.
The one-cell rows can then be filled in $4$, $6$, or $4$ ways,
respectively.  This gives
\[
2(4+6+4)=28.
\]

\item They may be in different long rows.  There are $2$ ways to
decide which long row contains $3$.  If both companions are $1$, the
one-cell rows can be filled in $4$ ways; if both are $2$, there are
$4$ ways; and if the companions are different, there are $2$ choices
for their assignment and $6$ fillings of the one-cell rows.  This gives
\[
2(4+4+2\cdot6)=40.
\]
\end{enumerate}
Thus this case contributes $28+40=68$.
\end{enumerate}

Adding all three cases gives
\[
\boxed{M_{(2^2,1^4),(3^2,1^2)}(h,m)=120+224+68=412.}
\]

\item \textbf{Solution.}
We count weakly increasing row fillings of shape $\lambda$ and content
$1^n$.  Since the entries $1,2,\ldots,n$ are all distinct, choose
which $\lambda_1$ entries go in the first row, which $\lambda_2$ of
the remaining entries go in the second row, and so forth.  After the
entries of a row have been chosen, their weakly increasing order is
forced; because the entries are distinct, it is actually strictly
increasing.  Thus
\[
\boxed{M_{\lambda,1^n}(h,m)
=\frac{n!}{\lambda_1!\cdots\lambda_\ell!}.}
\]

\item \textbf{Solution.}
The content $(n)$ means that every cell contains $1$.  This is allowed
because a row only has to be weakly increasing, so repeated entries
are permitted.  Every row is therefore forced, and there is exactly
one tableau of every shape $\lambda\vdash n$.  Thus
\[
\boxed{M_{\lambda,n}(h,m)=1}
\qquad\text{for every }\lambda\vdash n.
\]

\item \textbf{Solution.}
The content $(n-1,1)$ consists of $n-1$ copies of $1$ and one copy of
$2$.  Choose the row that will contain the unique $2$.  In that row,
weak increase forces the $2$ to be the final entry, with all earlier
entries equal to $1$.  Every other row consists entirely of $1$'s.

Thus each row gives exactly one tableau, depending on whether that row
is chosen to contain the $2$.  Since $\lambda$ has $\ell(\lambda)$
rows, including rows of equal length in different positions, we obtain
\[
\boxed{M_{\lambda,(n-1,1)}(h,m)=\ell(\lambda),}
\]
where $\ell(\lambda)$ is the number of parts of $\lambda$.

\item \textbf{Solution.}
Write
\[
p_\lambda=p_{\lambda_1}\cdots p_{\lambda_\ell},
\qquad
p_{\lambda_j}=\sum_{i\geq1}x_i^{\lambda_j}.
\]
For each row $j$, choosing the term $x_i^{\lambda_j}$ from
$p_{\lambda_j}$ means placing the entry $i$ in all $\lambda_j$ cells
of that row.  Therefore the relevant tableaux are exactly the fillings
in which \emph{every row is constant}.  There is no condition relating
one row to another.

After making one choice for every row, the exponent of $x_i$ is the
total number of cells filled with $i$.  Hence the resulting monomial
has exponent sequence $\mu$ exactly when the filling has content
$\mu$.  We have therefore proved the combinatorial interpretation
\[
\boxed{M_{\lambda,\mu}(p,m)
=\#\{\text{constant-row fillings of shape $\lambda$ and content $\mu$}\}.}
\]

\item \textbf{Solution.}
We count constant-row tableaux of shape $(2,2,1,1,1,1)$ and content
$(3,3,1,1)$.  Thus we need three $1$'s, three $2$'s, one $3$, and one
$4$.

The two long rows cannot be filled with $3$ or $4$, because a constant
row of length $2$ would use two copies of that entry, but the content
provides only one copy of each.  They also cannot both be filled with
$1$, since that would require four $1$'s, and they cannot both be
filled with $2$.  Therefore one long row must be $[1,1]$ and the other
must be $[2,2]$.

There are $2$ ways to decide which long row gets the $1$'s.  After
that choice, the four one-cell rows must contain one each of
$1,2,3,4$, and these four entries can be assigned to the four rows in
$4!$ ways.  Hence
\[
\boxed{M_{(2^2,1^4),(3^2,1^2)}(p,m)=2\cdot4!=48.}
\]

\item \textbf{Solution.}
The content $1^n$ uses each of the entries $1,2,\ldots,n$ exactly
once.  A constant row of length greater than $1$ would repeat its
entry, so such a row is impossible.  Consequently every row must have
length $1$, which means $\lambda=1^n$.

For this shape, assign the $n$ distinct entries to the $n$ one-cell
rows.  The first row has $n$ choices, the second has $n-1$, and so on,
for a total of $n!$ tableaux.  Therefore
\[
\boxed{M_{\lambda,1^n}(p,m)=
\begin{cases}
n!,&\lambda=1^n,\\
0,&\text{otherwise}.
\end{cases}}
\]

\item \textbf{Solution.}
The content $(n)$ means that all $n$ cells contain $1$.  Every row is
then constant automatically, and there are no choices to make.  Thus
there is exactly one such tableau for every shape $\lambda\vdash n$,
and
\[
\boxed{M_{\lambda,n}(p,m)=1}
\qquad\text{for every }\lambda\vdash n.
\]

\item \textbf{Solution.}
The content $(n-1,1)$ consists of $n-1$ copies of $1$ and one copy of
$2$.  Because rows must be constant, the row containing $2$ cannot
have length greater than $1$: otherwise it would contain at least two
$2$'s.  Thus the unique $2$ must occupy a one-cell row.  After choosing
that row, every other row is forced to contain only $1$'s.

There is one choice for each part of $\lambda$ equal to $1$.  If
$m_1(\lambda)$ denotes the number of such parts, then
\[
\boxed{M_{\lambda,(n-1,1)}(p,m)=m_1(\lambda),}
\]
where $m_1(\lambda)$ denotes the multiplicity of the part $1$ in
$\lambda$.

\end{enumerate}
\end{document}
